\subsection*{IF-13. Rational test vectors span the projected image of the fixed theta chain}

This statement addresses a different issue from the integral
coefficient transport in IF-08: whether a nonzero complex period on
the actual, fixed global theta chain can be detected by the rational
test vectors used in the source. It does not require the entries of
the row theta map to be algebraic. The central finite-group projector
must, however, be retained.

\begin{proposition}[IF-13: projected rational detection]
Let $K\hookrightarrow\mathbb C$ be a field embedding $\sigma$,
$V_K$ a finite-dimensional $K$-space, and
$j:V_K\longrightarrow H$ a $\sigma$-semilinear map whose image
spans the complex vector space $H$. Let $G$ be finite, let
$r:G\longrightarrow\operatorname{GL}_K(V_K)$ be a representation,
and let $a\in K[G]$. Write $q=r(a)$ and
$V_K^a=q(V_K)$. Let $\Phi:H\longrightarrow W$ be complex linear.
Suppose complex-linear operators $A_g$ on $W$ satisfy
\[
       \Phi(j(r(g)x))=A_g\Phi(j(x))
       \qquad(g\in G,\ x\in V_K).                         \tag{IF-13a}
\]
Set $Q=\sum_g\sigma(a_g)A_g$. Suppose $v\in\Phi(H)$ and $Qv=v$.
Then
\[
 v\in\operatorname{span}_{\mathbb C}
       \{\Phi(j(x)):x\in V_K^a\}.                         \tag{IF-13b}
\]
In particular, if a complex linear or conjugate-linear functional
$\ell$ has $\ell(v)\ne0$, then $\ell(\Phi(j(x)))\ne0$ for
some $x\in V_K^a$.

Apply this to a source realization $W_0$, a fixed family
$f:\texttt{FamCh}\ W_0\ u$ satisfying \texttt{RowOK} and
\texttt{LinkOnto}, and the model, character and reduction hypotheses
used in \texttt{PkHyp}. Put $R=\texttt{pk}(W_0)$ and
$A=W_0.P.A3.A$. Take $t=\mu_0$,
$K=D.\texttt{what}_2.\texttt{LL}$, $\sigma=\texttt{embLL}$,
$V_K=(A.\texttt{fam.R}\ t).\texttt{S.Forms}$,
$j=R.\texttt{jX}\ t$, and $\Phi=R.\texttt{thDn}\ t$.
Use exactly the source group-algebra element
\[
 a=e_X e_{P,\psi},\qquad
 e_{P,\psi}=\frac1{|P|}\sum_{h\in P}\psi(h)^{-1}h,
 \qquad
 e_X=\texttt{idemOf}(A.X),                                \tag{IF-13c}
\]
where $(P,\psi)=(W_0.Q.F.P1,W_0.Q.F.\psi1)$.
If
\[
 T\in\texttt{ThetaLiftSpaceZero}(W_0),\qquad
             QT=T,                                      \tag{IF-13d}
\]
then
\[
 T\in\operatorname{span}_{\mathbb C}
     \{\texttt{rowAt}(W_0,\mu_0,F):F\in K^{n_V}\}.
                                                            \tag{IF-13e}
\]
The fixed-point condition in (IF-13d) holds in particular if the
multilinear form corresponding to $T$ lies in the source
$R.\texttt{EEat}(\texttt{lvP}\ W_0)$: that space includes both
the $P,\psi$ level and the central $e_X$ cut. Thus every nonzero
linear period on such a target is detected by a rational test $F$.
The same conclusion holds with $F$ integral after multiplying it
by a nonzero integer; no bound on that clearing integer is asserted.
\end{proposition}

\begin{proof}
Choose $y\in H$ with $\Phi(y)=v$. Since the image of $j$ spans,
there are finitely many $c_i\in\mathbb C$ and $x_i\in V_K$ with
$y=\sum_i c_i j(x_i)$. Equation (IF-13a), semilinearity of $j$,
and linearity of $\Phi$ give, without any algebraicity assumption
on $\Phi$,
\[
 \Phi(j(qx_i))
  =\sum_g\sigma(a_g)\Phi(j(r(g)x_i))
  =Q\Phi(j(x_i)).                                        \tag{IF-13f}
\]
Consequently
\[
 v=Qv=\sum_i c_i\Phi(j(qx_i)),
\]
and every $qx_i$ lies in $V_K^a$. This proves (IF-13b). If all
the displayed test values under $\ell$ were zero, their linear
combination, with coefficients $c_i$ or $\overline{c_i}$ as
appropriate, would give $\ell(v)=0$, a contradiction. This also
proves the assertion for conjugate-linear functionals.

We now verify every identification needed for the source
specialization. The source results below are used with their
displayed hypotheses; this argument does not construct an
unconditional \texttt{RealPack}, \texttt{RowOK} or
\texttt{LinkOnto} witness.

First, the rational image $R.\texttt{jX}\ t$ spans
$(R.\texttt{HX}\ t).H$ under \texttt{IsModels}. This is the
untwisted content of \texttt{RealDat.kFormR\_of},
\texttt{Formalism/LocalToGlobal/IndefiniteReal/LinksLemmas.lean},
lines 148--154. Its proof uses the actual finite class-form
presentation, \texttt{ClassPres.extBij\_holds}, the model
embedding \texttt{HX.e}, and \texttt{models\_jX}. In particular
this is a full rational form of the definite top coefficient
space, before imposing either a theta-adjoint-image condition or a
finite-group projector. The conjugation twist used to define
\texttt{toCx} is irrelevant to this untwisted spanning assertion;
the proof constructs it before applying \texttt{Tw.span\_top'}.

Second, the actual fixed composite is the required $\Phi$.
The definitions of \texttt{Reach.step} and \texttt{LinkOnto}
refer to composition with the chosen row operators; they do not
assert only that variable theta kernels collectively span a target.
The source assumptions $m\ge2$ and $m$ even supply an integer
$n\ge0$ with $2n+2=m$. Given $T$ in (IF-13d),
\texttt{LinkOnto} supplies one
$G_0:(f.t\ \mu_0).\texttt{CT}$ with
$\texttt{down}(f,G_0)=T$. The \texttt{ioT} clause of
\texttt{RowOK.row} supplies membership of $G_0.1$ in the range of
\[
 R.\texttt{cX.toMap}\ R.\texttt{prX}\circ R.\texttt{HX.e}.
\]
Choose its preimage $y$. The supplied
\texttt{GlobalTheory.RowDescent.down\_eq},
\texttt{Development/AbelianVarietyTheory/Proofs/DecoratedLatticeProofs/}
\texttt{FamilyLinkDownward.lean}, lines 58--75, proves
\[
              R.\texttt{thDn}\ \mu_0\ y=T.                \tag{IF-13g}
\]
This is the full raw composite before the central projection in
$R.\texttt{Th}$. No inverse to a scalar period normalization, and
no rational preimage of $T$, has been used.

Third, (IF-13a) is exactly the supplied
\texttt{GlobalTheory.GoodFix.thDn\_act},
\texttt{EZeroThetaDescentEval.lean}, lines 63--77. On pure tensors
its right-hand operator is the source finite-group action
\texttt{actL} on the bottom row functions. Pure tensors span the
tensor product defining \texttt{SmT}; hence its pointwise
identity is an equality of the linear row maps. That theorem
uses \texttt{CuspRapidDecay}, the actual chain identities,
\texttt{IsModels}, and \texttt{grp\_red\_R}; it retains the
chosen Weil datum, level characters, measures and step vectors.
It therefore applies to the same $\Phi$ as (IF-13g).

Fourth, the actual input space of rational tests is the image of
$q=r(e_Xe_{P,\psi})$. The identification is
\texttt{CxPass.LevelOperators.goodR\_eq},
\texttt{LiftGoodPartComplexEqRange.lean}, lines 452 onward,
and the definition of
\texttt{BlockGood.goodR}; equivalently, the source's matrix
\texttt{vE} parametrizes that image. The exact map
\texttt{Face3.eV} is defined in
\texttt{ZeroPartEigenvalueConvention.lean}, lines 64--67, as
\[
 \texttt{basis.equivFun.symm}\circ
                 \texttt{vE.mulVecLin}.                   \tag{IF-13h}
\]
Thus every $qx_i$ occurring in the first part of the proof equals
$\texttt{eV}(F_i)$ for some $F_i\in K^{n_V}$. Although
\texttt{eV} has a dependent type involving \texttt{Zs}, its
formula contains no projection to \texttt{Zs}. Confusing it with
the separate map \texttt{toF} from an intersection with the
theta-adjoint image would incorrectly introduce a
\texttt{RatZ} hypothesis here.

Fifth, $q$-projected tests are centrally fixed. The members
in $A.X$ are the source's absolute irreducible characters, so
$e_X^2=e_X$ and $e_X$ is central. Hence $e_Xq=q$.
The exact source verification is \texttt{good\_fix},
\texttt{EZeroThetaDescentEval.lean}, lines 23--53. Applying
(IF-13a) termwise shows that
$R.\texttt{thDn}\ t\,(R.\texttt{jX}\ t\,\texttt{eV}(F_i))$
is fixed by the bottom central projector $e_0$.
The source's \texttt{chain\_Th},
\texttt{LinksChar.lean}, lines 289--290, says precisely that
$R.\texttt{Th}$ is the model of $e_0\circ R.\texttt{thDn}$.
Therefore
\[
 \Phi(j(\texttt{eV}(F_i)))
        =\texttt{rowAt}(W_0,\mu_0,F_i).                    \tag{IF-13i}
\]
The multilinear-to-linear conversion is
\texttt{lift\_postOn\_fixed},
\texttt{FamilyLinkDownward.lean}, lines 77--83; the final
\texttt{linkDown\_of\_rowOK} proof makes the same conversion.
Equations (IF-13b), (IF-13g), and (IF-13i) prove (IF-13e).

Finally verify the stated sufficient condition on $T$.
In \texttt{LinksDefs.lean}, lines 628--638, the source defines
$e_0$ using \texttt{idemOf}(A.X), and defines
$\texttt{EEat}(P)$ as the \texttt{goodPart} cut of its row
multiplicity space by this $e_0$. Thus $e_XT=T$. Its level
condition says exactly that the normalized finite-group action
of $h\in P$ is multiplication by $\sigma(\psi(h))$.
Indeed, a lift of $h$ contributes the inverse base-level character
in \texttt{actL}, while the extended level character contributes
the product of that base character and $\sigma(\psi(h))$;
the two base characters cancel. Averaging therefore gives
\[
 e_{P,\psi}T
   =\frac1{|P|}\sum_{h\in P}
        \sigma(\psi(h))^{-1}\sigma(\psi(h))T=T.
\]
Consequently $QT=e_Xe_{P,\psi}T=T$, as required.

For an integral test, only finitely many coordinates of the chosen
$F$ must be cleared in the number field $K$. Choose a positive
integer $N$ with $NF\in\mathcal O_K^{n_V}$. The map
\texttt{rowAt} is semilinear in $F$, and
$\sigma(N)=N\ne0$, so its nonzero period remains nonzero.
This last operation gives existence and does not give a unit
normalization or a valuation bound for $N$.
\end{proof}

\paragraph{Exact scope of the reduction.}
The unprojected \texttt{CXP} space has only the prescribed level
and character; it need not equal the span of \texttt{topOf}.
The formula for $e_X$ cannot be dropped. The proposition proves
that an actual nonzero bottom period with an anchor satisfying
(IF-13d) is detected by the exact rational source tests. It does
not construct such an anchor, the fixed-family realization, or
the $p$-adic factorization hypotheses for a divisibility estimate.
In particular a nonzero vector merely in
\texttt{ThetaLiftSpaceZero}, without its central cut, is not
enough for this argument. The proof also asserts no equality of
ordinary Lefschetz order indices with global period valuations.
